\section{Extended fiducial phase-space study}
\label{sec:fps}

This section reports a preliminary scalar central value, a three-prior dependence
envelope, and an extension region tested by a hidden-variable closure and
200 coverage toys. No corrected covariance is available for the extended
fiducial domain, so the result is limited to central values and these validation
tests.
The signal remains $\nu_\mu$ CC in the tracker interaction fiducial volume:
truth-muon kinematic cuts are removed, but the reconstructed MINOS-based
selection is unchanged. This is therefore an extended \emph{fiducial}
measurement domain, not unrestricted neutrino phase space.

The published phase space ($\theta_\mu<20^{\circ}$, $1.5<\pPar<\SI{60}{GeV/c}$,
$\pT<\SI{4.5}{GeV/c}$) hides a third of the signal: weighting the true
$\nu_\mu$~CC rate in the tracker fiducial volume, \SI{66.4}{\percent} lies
inside the cuts, \SI{22.3}{\percent} at $\pPar<\SI{1.5}{GeV/c}$,
\SI{11.3}{\percent} at $\theta_\mu>20^{\circ}$, and the high-$\pT$/$\pPar$
tails are negligible (Fig.~\ref{fig:fpsacc}). Because OmniFold's miss
treatment performs acceptance correction natively, removing the truth-level
kinematic cuts is a configurational change (an event-loop switch enlarging
the truth denominator from 32.8M to 49M events; the reconstruction-level
selection, including the MINOS match, is untouched), and the unfold
extrapolates into the unmeasured regions through the step-1 miss regressor.
The question is not \emph{whether} it can extrapolate but \emph{at
what prior cost}; this section quantifies that.

\begin{figure}[H]
  \centering
  \includegraphics[width=0.99\textwidth]{fps_acceptance_MEFHC}
  \caption{Extended-fiducial acceptance study (all twelve playlists):
  (a)~weighted truth rate (log$_{10}$), (b)~reconstruction efficiency, and
  (c)~the dead region with $\epsilon<\SI{2}{\percent}$, on the extended
  $(\pT,\pPar)$ grid (bin number; the catch bins are too wide for a linear
  axis). Cyan lines mark the published phase-space boundary. Efficiency is
  $\sim0.65$ and flat through the published region, degrades smoothly across
  the $\theta>20^{\circ}$ wedge, and dies in the $\pPar<\SI{1.5}{GeV/c}$
  strip and the high-$\pT$ corner; the dead region carries
  \SI{27.7}{\percent} of the total truth rate.}
  \label{fig:fpsacc}
\end{figure}

\paragraph{Validation battery.}
Three results establish the extension (Fig.~\ref{fig:fpscmp}):
\begin{enumerate}[itemsep=2pt]
  \item \textbf{Anchor.} Restricted to the published-phase-space sub-block,
        the extended-fiducial unfold reproduces a control unfold (standard
        omnifile, paper grid, identical settings) to $0.9994$ in integral,
        with a median per-cell ratio of $1.0013$ and median deviation
        \SI{\fpsAnchor}{\percent} over all 185 reported cells --- opening the
        phase space does not damage the measurement where it was already
        made. The control itself reproduces the frozen production total
        ($\SI{\sigTwoD}{cm^2/nucleon}$) exactly.
  \item \textbf{Closure.} A self-closure on the extended grid (simulated
        reconstruction as pseudo-data) returns recovered/truth $=1.0000$ in
        every cell including the extension --- no bookkeeping or
        normalization bias in the enlarged-denominator configuration. A
        \emph{hidden-variable} closure goes further: injecting a Gaussian
        truth bump in true $E_{\mathrm{avail}}$ (an axis the
        $(\pT,\pPar)$ unfold is blind to) and unfolding, the recovered
        $(\pT,\pPar)$ cross section reproduces the bump-reweighted truth to
        median \SI{0.17}{\percent} (90th percentile \SI{0.65}{\percent}) in
        the published cells and median \SI{0.77}{\percent} (max
        \SI{4.05}{\percent}) in the extension --- both well inside the
        tier-2 prior band.
  \item \textbf{Three-prior envelope.} Repeating the unfold with three
        priors --- MnvTune~v1, bare GENIE, and a NuWro-shaped prior (built
        by reweighting the truth $(\pT,\pPar)$ shape to two million raw
        NuWro events, which carry no phase-space cut) --- spans the total
        cross section by only $\pm\SI{1.5}{\percent}$. Per cell, the
        half-spread/mean is median \SI{2.9}{\percent} inside the published
        region but median \SI{7.9}{\percent} (90th percentile
        \SI{62}{\percent}) in the extension, concentrated in the
        zero-acceptance cells (Fig.~\ref{fig:fpsenv}).
  \item \textbf{Containment diagnostic.} With the estimator seed fixed, one 100/100 calibration--evaluation
        split, \SI{\covFPS}{\percent} of evaluation bin--toy cells contain the fixed
        truth using widths learned from the disjoint calibration half, compared with the
        \SI{\covNominal}{\percent} Gaussian nominal. Correlated bins and estimated
        calibration widths preclude a global binomial error or calibrated test; this probes the
        statistical closure-toy stream, not the full adopted band. Separately,
        the independent $C_{\mathrm{stat}}$ band contains \SI{\covFPSanalytic}{\percent}
        of closure bin--toy cells and is wider than the closure-toy spread. The
        same-ensemble \SI{68.9}{\percent} ($\langle|r|\rangle\approx0.79$) is retained as a
        pull/Gaussianity diagnostic. The published and extension regions have
        descriptive split-sample fractions of \SI{68.5}{\percent} and
        \SI{70.0}{\percent}, respectively; neither is used as a covariance-certification test.
\end{enumerate}

\begin{figure}[H]
  \centering
  \includegraphics[width=0.99\textwidth]{fps_pilot_compare_MEFHC}
  \caption{Extended-fiducial validation battery (all playlists):
  (a)~anchor ratio in the published sub-block (median $1.0013$),
  (b)~the MnvTune/bare-GENIE prior-swap ratio, and (c)~the extended-fiducial
  cross section (log$_{10}$). Panels (a) and (b) are \emph{ratios} (diverging
  scale centered at unity); (c) is an \emph{absolute} magnitude (sequential
  scale) --- hence the two color schemes. Black lines mark the published
  boundary; axes are bin number on the extended grid.}
  \label{fig:fpscmp}
\end{figure}

\begin{figure}[H]
  \centering
  \includegraphics[width=0.62\textwidth]{fps_prior_envelope_MEFHC}
  \caption{Three-prior dependence envelope (half-spread/mean per cell,
  in percent). Prior dependence is at the few-percent level wherever the
  detector has acceptance and reaches tens of percent only in the dead
  cells --- the quantitative basis for the two-tier reporting.}
  \label{fig:fpsenv}
\end{figure}

\paragraph{Result and reporting convention.}
The preliminary extended-fiducial CC-inclusive central value on the tracker
fiducial volume is
\begin{equation*}
  \sigma_{\mathrm{ext}} = \SI{\fpsSigma}{cm^2/nucleon}
  \quad\text{(central value)},
\end{equation*}
\SI{\fpsAbove}{\percent} above the restricted-phase-space
$\SI{\sigTwoD}{cm^2/nucleon}$. We adopt \textbf{two-tier reporting}:
\emph{tier~1} (acceptance-supported reporting) covers cells with $\epsilon\gtrsim2\%$ ---
the published region plus the measurable extension ($\sim\SI{6}{\percent}$
of the rate, chiefly the $\theta>20^{\circ}$ wedge and the
$\pPar\in[0.75,1.5]\,\si{GeV/c}$ edge); \emph{tier~2} (prior-extrapolated)
covers the dead cells ($\sim\SI{28}{\percent}$ of the rate), separately
flagged and carrying the three-prior envelope as an explicit additional band.
The summed central value includes both tiers; only tier~1 is interpreted as an
acceptance-supported measurement. No uncertainty magnitude is assigned beyond
the three-prior envelope because the extended-domain covariance has not been
constructed under the contract of \S\ref{sec:uq-combine} with
selection-complete lateral support.

\subsection{Recoil-PET feasibility study on the extended population}
\label{sec:fps-pet}
The PET feasibility study on the enlarged event population and reporting grid
remained recoil-only: step~1 saw non-muon reconstructed clusters and
step~2 saw truth hadrons with the muon removed. The muon coordinates used to
define and bin the extension were not classifier inputs. Those products are
therefore feasibility diagnostics for the recoil estimator;
their populated-bin counts, prior-transfer envelope and retraining floor are
not a final full-event measurement and are not used to argue that
low-efficiency cells are data-constrained.

A full-event estimator must satisfy two independent gates. It must first
carry a distinguished muon and validated event context through separate
observable reco/data and truth interfaces, and pass an omitted-variable stress
closure that shifts muon kinematics at fixed recoil. It must also consume the
canonical extended-fiducial event loops and $15\times19$ muon grid while
preserving the enlarged truth denominator, native misses and unchanged
reconstructed selection.

Its validation would need to repeat the ordinary and stress closures,
acceptance census, anchor, coverage and prior envelope before constructing
statistical, ML, vertical, retraining-response and selection-complete-lateral
uncertainties. Additional features can use more observed information, but
cannot create detector information in zero-efficiency cells.
